% context: 9th-grade physics classwork, Unit 1 measurement
% topic: rate problems — rate x time = quantity (and its rearrangements),
%        followed by a unit conversion; modeled on the 1-7Q early finisher problems
% structure: two pages, two problems per page
\documentclass[12pt, twoside]{article}
%\documentclass[12pt, twoside]{article}
\usepackage[letterpaper, margin=1in, headsep=0.2in]{geometry}
\setlength{\headheight}{0.6in}
%\usepackage[english]{babel}
\usepackage[utf8]{inputenc}
\usepackage{microtype}
\usepackage{amsmath}
\usepackage{amssymb}
%\usepackage{amsfonts}
\usepackage[nomessages]{fp} %\FPeval{\var-name}{2*sin(pi/6)}
\usepackage{siunitx} %units in math. eg 20\milli\meter
\usepackage{yhmath} % for arcs, overparenth command
\usepackage{tikz} %graphics
\usetikzlibrary{quotes, angles, arrows, arrows.meta}
\usepackage{graphicx} %consider setting \graphicspath{{images/}}
\usepackage{parskip} %no paragraph indent
\usepackage{enumitem}
\usepackage{multicol}
\usepackage{venndiagram}

\usepackage{fancyhdr}
\pagestyle{fancy}
\fancyhf{}
\renewcommand{\headrulewidth}{0pt} % disable the underline of the header
\raggedbottom
\hfuzz=2mm %suppresses overfull box warnings

\usepackage{hyperref}
\usepackage{float}

\title{Physics}
\author{Chris Huson}
\date{October 2026}

\fancyhead[LE]{\thepage}
\fancyhead[RO]{\thepage \\ Nickname: \hspace{2.25cm} \,\\ \,}
\fancyhead[LO]{La Scuola d'Italia / Huson / Physics \\* 9 October 2026}

\begin{document}

\subsubsection*{1.9 Rate Problems}

\textbf{Reference (provided).} \quad
rate $\times$ time $=$ quantity \quad (for motion: $r \cdot t = d$) \\[0.1cm]
1 km = 1000 m \quad 1 m = 100 cm \quad 1 h = 3600 s \quad 1 min = 60 s
\quad 1 gal $\approx$ 3.785 L

\vspace{0.2cm}

Solve each problem in two steps. First use rate $\times$ time $=$ quantity, or a
rearrangement of it. Then convert your result to the units the question asks for.
Round your final answer to the number of significant figures the data supports.

\vspace{0.3cm}

\begin{enumerate}

\item The Sun is $1.496 \times 10^{8}$ km from the Earth. Light travels at
  $3.00 \times 10^{8}$ m/s. How many minutes does sunlight take to reach the Earth?
  \vspace{7.5cm}

\item Hair grows about 15 centimeters per year. A student does not cut their hair from
  age 9 to age 15. By how many meters does it grow?

\newpage

\item A US kitchen faucet lets through at most 2.2 gallons of water per minute. Filling
  a large pasta pot takes 45 seconds. How many liters of water go into the pot?
  \vspace{9cm}

\item Usain Bolt's 100-meter world record is 9.58 seconds. What was his average speed in
  kilometers per hour?

\end{enumerate}
\end{document}
